% Generated draft: requires mathematical and source review.
\documentclass[11pt]{article}
\usepackage[margin=25mm]{geometry}
\usepackage{fontspec}
\setmainfont{Latin Modern Roman}
\usepackage{amsmath,amssymb,mathtools,unicode-math}
\setmathfont{Latin Modern Math}
\usepackage{enumitem,needspace,xurl,hyperref}
\setlength{\parindent}{0pt}
\setlength{\parskip}{5pt}
\setlength{\emergencystretch}{3em}
\newcommand{\N}{\mathbb{N}}
\newcommand{\Z}{\mathbb{Z}}
\newcommand{\Q}{\mathbb{Q}}
\newcommand{\R}{\mathbb{R}}
\newcommand{\C}{\mathbb{C}}
\newcommand{\F}{\mathbb{F}}
\newcommand{\A}{\mathbb{A}}
\newcommand{\PP}{\mathbb{P}}
\newcommand{\E}{\mathbb{E}}
\newcommand{\eps}{\varepsilon}
\newcommand{\dd}{\,\mathrm{d}}
\begin{document}
\section*{UnsolvedMath ID: 30004504}
\textbf{Source problem code:} catalogue entry
\subsection{Definitions and mathematical statement}
For a positive integer \(n\), let \(\mu(n)\) be the Möbius function: \(\mu(1)=1\), \(\mu(n)=0\) if a prime square divides \(n\), and \(\mu(n)=(-1)^r\) if \(n\) is a product of \(r\) distinct primes. For \(u\in\R\), put \(e(u)=\exp(2\pi i u)\). Let \(k\in\Z_{\geq1}\) be fixed, and let \(H=H(X)\) be a positive real-valued function such that \(H(X)\to\infty\) as \(X\to\infty\). Define
\[
M_{k,H}(x;X)=\sup_{\substack{P\in\R[t]\\ \deg P=k}}\left|\sum_{\substack{n\in\Z_{\geq1}\\ x\leq n\leq x+H(X)}}\mu(n)e(P(n))\right|.
\]
The summation includes both endpoints whenever they are integers. The supremum is over all real-coefficient polynomials of degree exactly \(k\), and the polynomial may vary with \(x\) and \(X\). No upper bound on \(H(X)\) is stated in the supplied record.

The supplied question asks whether, for every such fixed \(k\) and every such function \(H\), one has \(M_{k,H}(x;X)=o(H(X))\) for almost all \(x\in[X,2X]\).

The record does not specify whether \(x\) ranges over real or integer starting points, or explicitly define ``almost all.'' Under the interpretation of real starting points and an exceptional set of relative Lebesgue measure tending to zero, the precise question is whether, for every \(k\in\Z_{\geq1}\), every positive \(H(X)\to\infty\), and every \(\varepsilon>0\),
\[
\lim_{X\to\infty}\frac{1}{X}\operatorname{meas}\bigl\{x\in[X,2X]:M_{k,H}(x;X)>\varepsilon H(X)\bigr\}=0,
\]
where \(\operatorname{meas}\) denotes Lebesgue measure. If integer starting points are intended instead, the corresponding question replaces this displayed limit by
\[
\lim_{X\to\infty}\frac{1}{X}\#\bigl\{x\in[X,2X]\cap\Z:M_{k,H}(x;X)>\varepsilon H(X)\bigr\}=0.
\]
These are explicit interpretations requiring source confirmation; the supplied record does not select between them. The limits are taken separately for each fixed \(k\) and each admissible \(H\), with no asserted uniformity over these choices.
\subsection{Short English statement}
Does the Möbius function exhibit cancellation, uniformly over real polynomial phases of each fixed degree \(k\geq1\), on almost every interval \([x,x+H(X)]\) with \(x\in[X,2X]\), for every length function \(H(X)\to\infty\), however slowly? The meaning of almost every starting point requires confirmation. The supplied record labels the problem open; this draft does not independently verify that status.
\subsection{Sources}
\begin{itemize}
\item Local UnsolvedMath registry: \texttt{data/unsolvedmath/problems.json}, permanent ID \(30004504\), source problem code \texttt{catalogue entry}, title ``Short-Interval Möbius Orthogonality to Polynomial Phases.'' The supplied collection is ``Oberwolfach Reports Open Problems.'' No individual proposer is named; the supplied proposed year is 2020.
\item Supplied source citation: \emph{Automorphic Forms and Arithmetic (2021)}, Oberwolfach Reports, pp.\ 1317--1383, DOI 10.4171/owr/2020/26. \url{https://doi.org/10.4171/owr/2020/26}. The background describes Oberwolfach Reports as documenting mathematical workshops and their open-problem discussions. The supplied literature list also labels this reference ``Oberwolfach Report 26/2020.''
\item Supplied literature reference: \emph{Higher uniformity of bounded multiplicative functions}. \url{https://arxiv.org/abs/2007.15644}.
\item Supplied literature reference: \emph{Higher uniformity of arithmetic functions II}. \url{https://arxiv.org/abs/2411.05770}.
\item Supplied literature assessment, checked at the exact supplied date \texttt{2026-08-22}: the status is recorded as open; the full short-interval polynomial-phase conjecture permits arbitrarily slowly growing intervals, while the cited primary work is described as proving only power-scale ranges, with a logarithmic-scale barrier remaining. This is supplied dated curation derived from the OWR record, not an independently verified literature assessment. Time-sensitive status and the cited mathematical claims require independent verification.
\end{itemize}
\end{document}
